\subsection{Wavelet Tree}

Wavelet Tree is a Data Structure mainly for finding the $k$'th minimum value in some range $[l, r]$. The tree is built dynamically, meaning that we will use pointers.

\begin{lstlisting}
struct Wavelet {
    vector <int> p;
    Wavelet *l, *r;
};
\end{lstlisting}

When we build the tree, we check if the $b$'th bit is 0 or 1. If it is 0, we put that element into the left side. Otherwise, we put in the right side. The vector \texttt{p} is a prefix sum array counting number of elements in the left side.

\begin{lstlisting}
void build(const vector <int> &A, int b) {
    if (b == 0 || A.empty()) return;
    l = new Wavelet();
    r = new Wavelet();
    int c = 0;
    vector <int> L, R;
    for (int x : A) {
        if (x & (1 << (b-1))) R.emplace_back(x);
        else L.emplace_back(x), c++;
        p.emplace_back(c);
    }
    l->build(L, b-1);
    r->build(R, b-1);
}
\end{lstlisting}

When we query, we count the number of elements in the left side in range $l$ to $r$. If the number of elements in the left side is more than or equal to $k$, we recursive to the right side. Otherwise we go to the right side. We have to change the value of $l, r, k$ too when moving down to the right side.

\begin{lstlisting}
int qr(int l, int r, int k, int b) { // [l, r), k, 0-index
    if (b == 0 || p.empty()) return 0;
    int pr = (r > 0 ? p[r-1] : 0);
    int pl = (l > 0 ? p[l-1] : 0);
    int cnt = pr - pl;
    if (k < cnt) return this->l->qr(pl, pr, k, b-1);
    else return (1 << (b-1)) + this->r->qr(l-pl, r-pr, k-cnt, b-1);
}
\end{lstlisting}

Note that these functions are inside the \texttt{Wavelet} struct.

\break